\documentclass{article}
\usepackage[utf8]{inputenc}
\usepackage{graphicx}
\usepackage{tikz}
\usetikzlibrary{calc}
\usepackage{graphicx}
\usepackage{xcolor}
\usepackage{tcolorbox}
\title{\textbf{}}


\begin{document}

\begin{tikzpicture}
[remember picture, overlay] \draw[line width = 4pt] ($(current page.north west) + (0.5in,-0.5in)$) rectangle ($(current page.south east) + (-0.5in,1in)$);
\end{tikzpicture}

\begin{center}
    
    {\LARGE \bf \textcolor{black}{RSA Cryptography}\par}

    \vspace{2cm}

    {\large \em Submitted by }\\
    \vspace{0.5cm}
   
   \begin{tcolorbox}
   \begin{center}
    \textbf{Arnish Satasiya (202001031)\\
        Jay Kuvadiya (202001042)\\
        Meet Sutariya (202001046)\\
        Deep Kanani (202001454)}
    \end{center}

    
    \end{tcolorbox}
    
    \vspace{1cm}
    {\large \em Submitted to }\\ 
    \vspace{0.5cm}
    {\LARGE \bf Prof. Manish K. Gupta}\\
    \vspace{0.5cm}
    {\LARGE \bf Prof. Rahul Muthu}\\
    \vspace{1cm}
    {\large \bf Course: Discrete Mathematics}\\ 
    \vspace{0.5cm}
    {\large \bf Course code: SC205}\\
    \vspace{1cm}
    \begin{center}
     \includegraphics[width=100pt]{Daiict.png} 
    \end{center}
    
\end{center}
\pagebreak 
\begin{tikzpicture}
[remember picture, overlay] \draw[line width = 4pt] ($(current page.north west) + (0.5in,-0.5in)$) rectangle ($(current page.south east) + (-0.5in,1in)$);
\end{tikzpicture}
\section{Introduction\\}
	
\begin{itemize}
\item RSA is used for transmission of data (messages) securely . 
\item The RSA algorithm, is an algorithm for Public Key cryptography .
\item In public-key cryptography, users reveal a public encryption key so that other users in the system are able to send private messages to them, but each user has their own private decryption key . 
\item The key to ensuring privacy in a public-key cryptosystem is for it to be extremely difficult to derive the decryption key from the publicly available encryption key . 
\item The algorithm works by exploiting concepts from Number Theory ,\\ including \textbf{Fermat’s Little Theorem , Sieve Algorithm and \\Extended Euclid Algorithm}   
\end{itemize}
\vspace{1cm}
\section{Mathematical Problem\\}

    Nowadays Countries or any company need a way to communicate without anyone except receiver being able to decipher what exactly they are communicating. How the communication worked ?
    
    The ans is that they started out by establishing a cipher that both parties would use, and from there one ally would use the numbers correlating to a certain number to create a message. Once the message was received from the other ally that group would use the cipher established to decode the message.

\pagebreak

\begin{tikzpicture}
[remember picture, overlay] \draw[line width = 4pt] ($(current page.north west) + (0.5in,-0.5in)$) rectangle ($(current page.south east) + (-0.5in,1in)$);
\end{tikzpicture}
\section{Solution\\}

\begin{itemize}
\item Here, we have try to solve this problem using our Discrete Mathematics knowledge.
\item Using the concept of RSA cryptography we can solve this real world problem.\\
\end{itemize}


{\LARGE \bf \textcolor{black}{How RSA works ?}\\}


\begin{itemize}
\item RSA works on two different keys i.e. Public Key and Private Key. 
\item The Public Key is accessible by everyone and Private key is kept private. 
\item Here, we describe two methods to perform RSA .   \\
\end{itemize}

\subsection{ Method 1 :  \\}

\begin{itemize}
\item If you want to send the message to one particular person , For that you have to encrypt your message using that person's public key n and e .
\end{itemize}

 \{ 
	Here, n is multiplication of two Prime numbers (p and q)  and e must 	be\\        selected by satisfying this condition ---\(>\) 
	( gcd of (p - 1)*(q - 1) and e  = 1 ) 
\}
\begin{itemize}
\item After that , encrypted message will be sent to that particular person and that person will decrypt the message using his/her private key d . Person will get the exact message that you have sent.
\end{itemize}

\begin{itemize}
\item By using this method the message will be shown to the only one person whom you want to send the message.
\end{itemize}

\begin{itemize}
\item we can understand this method by the example given below , In this example... Alice want to send a number 896524 to bob .
\end{itemize}

\begin{itemize}
\item for that Alice have used the public key n and e of Bob and encrypted the message , encrypted message 36982875 is sent to Bob.
\end{itemize}

\begin{itemize}
\item Now , Bob will decrypt the message using his private key d and he'll get the original message 896524 that Alice has sent.
\end{itemize}
\pagebreak

\begin{tikzpicture}
[remember picture, overlay] \draw[line width = 4pt] ($(current page.north west) + (0.5in,-0.5in)$) rectangle ($(current page.south east) + (-0.5in,1in)$);
\end{tikzpicture}

\begin{center}
\includegraphics[width=300pt]{RSA.png} 
\end{center}

\vspace{0.5cm}

\begin{itemize}
\item To encrypt a message using the RSA algorithm, the general method of the encryption procedure is as follows: 
\end{itemize}

\begin{tcolorbox}

\textbf{Simple RSA}\\\\
 First of all select message (\textcolor{red}{\textbf{M}}) that you ( here Alice ) want to send someone (here Bob).\\
 To encrypt the message , compute \( M^e \) mod n .\\\\

                 Cipher text \textcolor{blue}{\textbf{C}} = \textcolor{red}{\( M^e \)} mod n    \hspace{0.5cm} ( where n and e is Bob's Public Key )\\\\

     n = p*q    ( p and q are large Prime numbers ) \\
    \textbf{ \textcolor{blue}{gcd}} ( n , ( p-1 )*( q-1 ) ) = 1 .\\
    
    \end{tcolorbox}
    
    --\(>\) Now , C will be sent to Bob .
    
    \pagebreak
    
    \begin{tikzpicture}
[remember picture, overlay] \draw[line width = 4pt] ($(current page.north west) + (0.5in,-0.5in)$) rectangle ($(current page.south east) + (-0.5in,1in)$);
\end{tikzpicture}

    \begin{tcolorbox}
    
    To decrypt the C ( encrypted message ) Bob will do calculation that is given below :\\\\

                                  D = ( \(C^d\) ) mod n      \hspace{1cm} ( where d is Bob's Private Key )\\\\
where D is decrypted message .\\\\

      d$\cdot $e$ \equiv $ 1\hspace{0.25cm} (mod n) \\\\

So, by this Bob will get message that Alice want to sent him .\\

  \end{tcolorbox}

Code (in c++) : ( given in RSA1.cpp )\\\\
\centering\textbf{{\LARGE \bf \textcolor{black}{Limitations of this Code :}\\}}
\vspace{1cm}
\begin{tcolorbox}



\begin{itemize}
\item This code is valid only if we want to encrypt integers , but we can't use this for the characters .
\end{itemize}

\begin{itemize}
\item We can encrypt or decrypt the number in a limited range , if the number is to large (more than 8 digits) this code we'll not work .
\end{itemize}

\begin{itemize}
\item In the real world scenario , RSA only based on numbers won't so helpful so this is the main limitations of this code.\\
\end{itemize}

  \end{tcolorbox}
\pagebreak
    \begin{tikzpicture}
[remember picture, overlay] \draw[line width = 4pt] ($(current page.north west) + (0.5in,-0.5in)$) rectangle ($(current page.south east) + (-0.5in,1in)$);
\end{tikzpicture}
\centering\textbf{{\LARGE \bf \textcolor{black}{Updated RSA}\\}}
\vspace{1cm}
\begin{tcolorbox}

\begin{itemize}
\item  For solving the limitations of Simple RSA code, we have written this code.
\end{itemize}

\begin{itemize}
\item In this code , We have take input as string . By using string , we can also encrypt and decrypt very large length of message .
\end{itemize}

\begin{itemize}
\item we can also encrypt or decrypt the characters , symbols , numbers etc.
\end{itemize}

\begin{itemize}
\item Like the previous code there are same two methods for encrypt the message , This code is very similar to the previous code .
\end{itemize}

\begin{itemize}
\item In this code , we'll encrypt the message by their ASCII value , we'll first make 2-2 pairs of the characters in the message , then we'll write their ASCII value and we'll take the public key n and e of the person whom you want to send the message for encrypting. 
\end{itemize}
\begin{itemize}
\item For that first you should understand how we have used vector , we'll push back all the pairs in the vector .
\end{itemize}

\begin{itemize}
\item If the length of the string is odd then the last pair won't be made so  we'll push back only that perticular character's ASCII value in the vector.
\end{itemize}

\begin{itemize}
\item Now , one by one we'll encrypt every numbers of the vector using n and e as shown in previous code.
\end{itemize}

\begin{itemize}
\item After encrypting the message , we will print all the elements in the vectors by ASCII values converting to the characters . By that we'll get the encrypted message.
\end{itemize}

\begin{itemize}
\item Now for decrypt this encrypted message the receiver have to insert the decryption private key d .
\end{itemize}

\begin{itemize}
\item By Inserting the correct decryption key , Message will be decrypted successfully and the only person whom you have sent the message able to see the correct message .\\
\end{itemize}

\end{tcolorbox}

Code ( in c++ ) : ( refer RSAfinal.cpp )
\pagebreak

 \begin{tikzpicture}
[remember picture, overlay] \draw[line width = 4pt] ($(current page.north west) + (0.5in,-0.5in)$) rectangle ($(current page.south east) + (-0.5in,1in)$);
\end{tikzpicture}
\centering\textbf{{\LARGE \bf \textcolor{black}{Some Algorithms used in RSA}\\}}
\vspace{1cm}
\centering\textbf{{\LARGE \bf \textcolor{black}{1. Sieve Algorithm}\\}}
\vspace{1cm}
\begin{itemize}
\item We can get Prime numbers ( p \& q ) from most powerful Sieve Algorithms :  ( refer Sieve.cpp ) 
\end{itemize}

\begin{tcolorbox}
\begin{itemize}
\item The sieve of Eratosthenes is one of the most Powerful Algorithm to find all prime numbers smaller than N, even when N is very large number.\\
\end{itemize}
\end{tcolorbox}
\vspace{1cm}

\centering\textbf{{\LARGE \bf \textcolor{black}{2. power function}\\}}
\vspace{1cm}

\begin{itemize}
\item For understanding of pow() function \\: ( refer power.cpp )
\end{itemize}

\begin{tcolorbox}

\begin{itemize}
\item This algorithm is use for find the modulo of big numbers .
\item ( \(N^M\) modulo P) , where N, M and P are large numbers .\\
\end{itemize}
\end{tcolorbox}
\pagebreak
    \begin{tikzpicture}
[remember picture, overlay] \draw[line width = 4pt] ($(current page.north west) + (0.5in,-0.5in)$) rectangle ($(current page.south east) + (-0.5in,1in)$);
\end{tikzpicture}

\centering\textbf{{\LARGE \bf \textcolor{black}{3. Extended Euclidean Algorithm}\\}}
\vspace{1cm}              


\begin{itemize}
\item For Find Modular Multiplicative Inverse ( for finding d )\\ : ( refer MMI.cpp )\\
\end{itemize}

\begin{tcolorbox}

\begin{itemize}
\item We have two integers a and m and we want to find inverse of a mod m .\\
\end{itemize}

\begin{itemize}
\item The multiplicative inverse of a modulo m exists if and only if a and m are relatively prime ( gcd ( a, m ) = 1 )\\
\end{itemize}

and satisfy this equation : a$\cdot $x$ \equiv $ 1\hspace{0.25cm} ( mod m )

\begin{itemize}
\item \{ where , x is inverse of a mod m \}
\end{itemize}

         
\begin{itemize}
\item x should be in range \{ 1, 2, .....upto m-1 \}
\end{itemize}

\begin{itemize}
\item By Extended Euclidean Algorithm, For any two integers A and B, We can find integers s and t, such that:

                                                s·A + t·B = gcd( A, B )  
\end{itemize}
      
( s and t are know as  Bezout Coefficients )\\

   s$\cdot $a + t$\cdot $m  = 1  \hspace{1cm} \{ gcd ( a , m ) = 1 \}\\                            
  ( s$\cdot $a + t$\cdot $m ) mod m = 1 mod m\\
   s$\cdot $a mod m = 1 mod m\\                 
   a$\cdot $s $ \equiv $ 1 mod m\\                            

So, inverse of \textbf{a mod m} is\textbf{ s.}
\end{tcolorbox}


  
\pagebreak














\end{document}
